\section*{Bab \ref{ch:g11:forces-and-motion} --- Gaya dan Gerak}

\begin{solution}{exo:g11:forces-and-motion:1}
(a) Kecepatannya tetap: \omterm{def:g11:forces-and-motion:netforce}{gaya totalnya} nol. (b) Lajunya berubah: bukan
nol, berlawanan dengan geraknya. (c) Arahnya berubah (kecepatan adalah
vektor): bukan nol, menuju sisi dalam tikungannya. (d) Kecepatannya
tetap: nol --- \omterm{def:g10:universal-gravitation:weight}{berat} dan hambatan udaranya \omterm{def:g10:inertia:compensate}{saling meniadakan}.
\end{solution}

\begin{solution}{exo:g11:forces-and-motion:2}
$d_1/d_2 = 1.0/3.0$ dengan $d_1 + d_2 = \qty{0.80}{m}$:
$d_1 = \qty{0.20}{m}$ dari bola \qty{3.0}{kg}. Jika massanya sama:
tepat di tengah, \qty{0.40}{m} dari masing-masing.
\end{solution}

\begin{solution}{exo:g11:forces-and-motion:3}
Dengan $F\Delta t$ yang sama dan massa tiga kali lipat: $\Delta v$-nya
tiga kali lebih kecil ketika bermuatan. Untuk menyamai $\Delta v$ troli
yang kosong, doronglah tiga kali lebih kuat (atau tiga kali lebih lama).
\end{solution}

\begin{solution}{exo:g11:forces-and-motion:4}
$\Delta \vect v$: \qty{2}{m/s}, terarah \emph{ke belakang} (dari
\num{10} menjadi \qty{8}{m/s}). \omterm{def:g11:forces-and-motion:netforce}{Gaya totalnya} menunjuk ke arah yang
sama: berlawanan dengan $\vect v$ --- kasus 2, yaitu melambat.
\end{solution}

\begin{solution}{exo:g11:forces-and-motion:5}
(a) Statis, berlawanan dengan doronganmu (belum ada peluncuran).
(b) Kinetis, berlawanan dengan peluncurannya. (c) Statis, menunjuk
\emph{ke atas} lerengnya --- melawan peluncuran yang sedang diusahakan
gravitasi.
\end{solution}

\begin{solution}{exo:g11:forces-and-motion:6}
$\Delta v \propto F\Delta t/m$: (a) \qty{8.0}{m/s}; (b) \qty{2.0}{m/s};
(c) \qty{2.0}{m/s}; (d) $4.0 \times \frac{3 \times 2}{6} =
\qty{4.0}{m/s}$ --- ketiga perubahannya saling menghapus.
\end{solution}

\begin{solution}{exo:g11:forces-and-motion:7}
$v' = \sqrt{8.0^2 + 6.0^2} = \qty{10.0}{m/s}$, pada
$\tan\theta =
6.0/8.0$, yaitu $\theta \approx 37^\circ$ dari timur ke
arah utara. Kecepatan berjumlah sebagai \emph{vektor}: sumbangan yang
tegak lurus dipadukan menurut Pythagoras, bukan dengan menjumlahkan
besarnya.
\end{solution}

\begin{solution}{exo:g11:forces-and-motion:8}
Yang sudah tertetapkan sebelumnya: massa bolanya dan $\Delta v$-nya
(dari laju datangnya sampai nol), sehingga hasil kali $F\Delta t$ pun
tertetapkan. Memanjangkan $\Delta t$ dari \qty{0.02}{s} menjadi
\qty{0.08}{s} membagi \omterm{def:g10:inertia:force}{gaya} rata-ratanya dengan $4$.
\end{solution}

\begin{solution}{exo:g11:forces-and-motion:9}
\omterm{def:g11:forces-and-motion:friction}{Gaya gesek statis} \omterm{def:g10:relative-motion:trajectory}{lintasan} pada sepatunya: telapak kakinya mendorong
tanah ke belakang, dan \omterm{def:g10:inertia:contact}{gaya sentuh} yang timbal balik mendorong si pelari
ke depan. Paku sepatunya menaikkan nilai terbesar \omterm{def:g10:inertia:inventory}{gaya gesek} statisnya
(tidak tergelincir pada usaha penuh). Di atas es tanpa \omterm{def:g10:inertia:inventory}{gesekan} tidak ada
\omterm{def:g10:inertia:force}{gaya} luar mendatar sama sekali, sehingga kecepatan \omterm{def:g11:forces-and-motion:com}{pusat massanya} tidak
dapat berubah --- tidak ada \omterm{def:g10:orders-of-magnitude:prefixes}{awalan}, tidak ada langkah.
\end{solution}

\begin{solution}{exo:g11:forces-and-motion:10}
Titik sentuh roda yang menggelinding tidak meluncur: \omterm{def:g10:inertia:inventory}{gaya gesek} statis,
yang lebih kuat daripada \omterm{def:g11:forces-and-motion:friction}{gaya gesek kinetis} sebuah gelincir
(\cref{def:g11:forces-and-motion:friction}) --- perhentiannya lebih
pendek. \omterm{def:g11:forces-and-motion:friction}{Gaya gesek kinetis} ban yang tergelincir menunjuk berlawanan
dengan peluncurannya, berapa pun sudut rodanya: memutar kemudinya tidak
mengubah apa pun, dan kemudinya pun hilang.
\end{solution}

\begin{solution}{exo:g11:forces-and-motion:11}
Tegangan kawatnya, yang menunjuk ke dalam ke arah atletnya. Setelah
dilepaskan, bolanya terbang menyusuri garis lurus sepanjang garis
singgungnya, dengan laju yang dimilikinya (kelembaman). ``Kawat'' milik
Bulan adalah tarikan gravitasi Bumi; matikanlah gravitasinya dan Bulan
pergi menyusuri garis singgungnya, lurus dan beraturan, selamanya.
\end{solution}

\begin{solution}{exo:g11:forces-and-motion:12}
Dari ujung atasnya: $x = \dfrac{1.0 \times 0.60 + 2.0 \times 1.20}{3.0}
= \qty{1.00}{m}$. Jaraknya ke kedua pusat bagiannya: \qty{0.40}{m}
(tongkatnya) dan \qty{0.20}{m} (kepalanya), dengan nisbah $2 : 1$ ---
kebalikan nisbah massanya $1.0 : 2.0$, sebagaimana dituntut
\cref{def:g11:forces-and-motion:com}.
\end{solution}

\begin{solution}{exo:g11:forces-and-motion:13}
\emph{1.} Dengan $F$ dan $\Delta v$ yang sama serta $m$ dua kali lipat:
vannya memerlukan waktu dua kali lipat. \emph{2.} Dengan laju rata-rata
yang sama selama waktu dua kali lipat: jaraknya dua kali lipat.
\emph{3.} Dengan $\Delta v$ dua kali lipat pada $F$ dan $m$ yang sama:
waktunya dua kali lipat; tetapi laju rata-ratanya pun menjadi dua kali
lipat, sehingga jaraknya dikalikan $4$. Menggandakan lajumu
melipatempatkan jarak pengeremanmu.
\end{solution}

\begin{solution}{exo:g11:forces-and-motion:14}
\emph{1.} $T = 27.3 \times \qty{86400}{s} \approx \qty{2.36e6}{s}$,
sehingga $v = \dfrac{2\pi \times \num{3.84e8}}{\num{2.36e6}} \approx
\qty{1.0e3}{m/s} \approx \qty{1}{km/s}$. \emph{2.} Ke arah Bumi (ke
dalam); yaitu tarikan gravitasi Bumi. \emph{3.} Bulan memang
\emph{sedang} jatuh: perubahan kecepatannya menunjuk ke Bumi pada setiap
saat. Laju singgungnya semata-mata memastikan bahwa Bulan terus meleset
--- Bulan jatuh \emph{mengitari} Bumi dan bukan jatuh ke dalamnya.
\end{solution}

\begin{solution}{exo:g11:forces-and-motion:15}
Dengan $m$ dan $\Delta v$ yang sama, $F\Delta t$ pun tertetapkan:
$\Delta t$ yang lima puluh kali lebih lama di atas busa berarti \omterm{def:g10:inertia:force}{gaya}
rata-rata yang lima puluh kali lebih kecil. Menekuk lutut memanjangkan
$\Delta t$ pendaratannya dari sebuah sentakan menjadi lenturan yang
panjang; matras pesenam, secara mekanis, terbuat dari milisekon
tambahan.
\end{solution}

\begin{solution}{pb:g11:forces-and-motion:1}
\textbf{1.} $50/3.6 \approx \qty{13.9}{m/s}$ (sebutlah \qty{14}{m/s}).
\textbf{2.} $\Delta v \approx \qty{14}{m/s}$: penumpangnya bergerak
\qty{14}{m/s} sebelumnya dan nol sesudahnya --- kedua ujungnya
ditetapkan tabrakannya, sehingga tidak ada alat yang dapat menyentuh
$\Delta v$.
\textbf{3.} $\Delta v \propto F\Delta t/m$
(\cref{thm:g11:forces-and-motion:secondlaw}) dengan $\Delta v$ dan $m$
yang tertetapkan memaksa $F\Delta t$ tertetapkan pula. Satu-satunya
kenop bagi perancangnya adalah $\Delta t$.
\textbf{4.} $0.150/0.050 = 3$: zona remuknya membagi \omterm{def:g10:inertia:force}{gaya} rata-ratanya
dengan $3$.
\textbf{5.} Ke belakang, berlawanan dengan geraknya --- kecepatannya
menyusut dari \qty{14}{m/s} menjadi nol.
\textbf{6.} Bagian depan yang melipat adalah mesin untuk membuat
perhentiannya berlangsung lebih lama: logam yang remuk membelikan
milisekon.
\textbf{7.} Dengan $\Delta v$ dan $m$ yang sama serta $\Delta t$ tiga
kali lebih pendek: boneka di mobil kuno itu menerima \omterm{def:g10:inertia:force}{gaya} tiga kali
lipat.
\textbf{8.} Laju rata-ratanya $\approx \qty{6.9}{m/s}$; jaraknya
$\approx
6.9 \times 0.150 \approx \qty{1.0}{m}$ --- kira-kira sepanjang
satu \omterm{def:g10:orders-of-magnitude:unit}{meter} mobil yang dapat berubah bentuk yang memang disediakan zona
remuk yang sebenarnya.
\textbf{9.} Moncong sepanjang \qty{10}{m} tidak dapat dikemudikan dan
tidak dapat diparkir, dan massanya sendiri akan memperburuk setiap
tabrakan. Ruang penumpangnya tidak boleh remuk: ruang itu menjaga
ruang bertahan hidupnya.
\textbf{10.} ``\dots jadi kamu harus memilih $\Delta t$.''
\textbf{11.} $0.150/0.005 = 30$: perhentian pada kaca depannya tiga
puluh kali lebih kejam daripada perhentian yang bersabuk.
\textbf{12.} Sabuk yang sedikit meregang memanjangkan $\Delta t$ milik
penumpangnya sendiri sedikit lagi (sekaligus menyebarkan \omterm{def:g10:inertia:force}{gayanya}); tali
baja akan memaksakan perhentian terpendek milik badan mobilnya kepada
penumpang yang paling lunak.
\textbf{13.} $0.080/0.010 = 8$: kantong udaranya membagi \omterm{def:g10:inertia:force}{gaya} pada
kepalanya dengan $8$.
\textbf{14.} Dengan $\Delta v$ dan $\Delta t$ yang sama: \omterm{def:g10:inertia:force}{gayanya}
berskala dengan massanya, yaitu $70/20 = 3.5$ kali lebih besar bagi si
dewasa.
\textbf{15.} Perubahan yang dituntut: $14/0.150 \approx \qty{93}{m/s}$
tiap sekon --- sekitar $9.5$ kali yang dihasilkan gravitasi
(\qty{9.8}{m/s} tiap sekon). Lengannya harus menarik sekitar $9.5$ kali
\omterm{def:g10:universal-gravitation:weight}{berat} bayinya: seperti memangku beban \qty{95}{kg}. Tidak ada lengan
yang sanggup; bayinya terbang ke depan.
\textbf{16.} Sabuk orang dewasa menyalurkan \omterm{def:g10:inertia:force}{gayanya} ke pinggul dan
dada; pada seorang anak sabuk itu akan membebani perut dan lehernya.
Kursi anak mengerjakan \omterm{def:g10:inertia:force}{gayanya} yang (lebih kecil) pada bagian tubuh
yang kuat dan bantalannya memanjangkan $\Delta t$ lebih jauh lagi.
\textbf{17.} Jika ruang penumpangnya berubah bentuk, penumpangnya
kehilangan ruang bertahan hidupnya lalu terhantam oleh rangkanya
sendiri. Tidak ada apa pun yang boleh menghantam ruang yang kaku itu
secara langsung: perhentian penumpangnya ditangani sabuk dan kantong
udaranya ($\Delta t$ yang panjang), sementara ruang itu menambatkan
kedua ujung yang remuk.
\textbf{18.} Karena setangkup, kedua mobilnya berhenti pada bidang
sentuhnya: tiap pengemudinya mengalami $\Delta v \approx \qty{14}{m/s}$
sepanjang perhentian zona remuknya --- pada dasarnya sama dengan uji
dindingnya, bukan ``tabrakan yang dua kali lipat''.
\textbf{19.} \omterm{def:g10:inertia:force}{Gayanya} sama besar dan berlawanan, tetapi
$\Delta v \propto F\Delta t/m$: mobil yang kecil, dengan $m$ yang
kecil, menderita perubahan kecepatan yang besar.
\textbf{20.} Zona remuk, sabuk pengaman yang meregang, kantong udara,
bantalan kursi anak --- masing-masing memanjangkan besaran yang sama:
$\Delta t$.
\textbf{21.} Tabrakannya menetapkan massanya dan perubahan
kecepatannya, sehingga menetapkan pula hasil kali \omterm{def:g10:inertia:force}{gaya} $\times$ waktu:
tumbukan yang lembut adalah tumbukan yang lama --- panjangkanlah
waktunya, dan \omterm{def:g10:inertia:force}{gayanya} turun dengan nisbah yang sama.
\end{solution}
